\documentclass[11pt,a4paper]{article}

\usepackage{pifont}
\usepackage{amsmath}
\usepackage{amssymb}
\usepackage{psfig}
\usepackage{graphicx}
\usepackage{array}
\usepackage{fancyheadings}
\usepackage{here}
\usepackage{eepic,epic}
\usepackage[english]{babel}

\input{YRpersdf}

\oddsidemargin -0.9cm
\evensidemargin -0.9cm
\topmargin -1cm
\textheight 22.5cm
\textwidth 17.6cm
\headheight 1.0cm

% principal notations



\begin{document}

\begin{center}
  \includegraphics[width=10cm,angle=0]{getfem_logo.eps}\\[2cm]
  \fbox{\huge \sc Interpolation and Comparison program} \\[0.5cm]
  { \large Yves \sc Renard\footnote{ \it MIP, INSAT, Complexe scientifique de Rangueil, 31077 Toulouse, France, Yves.Renard@gmm.insa-tlse.fr } } \\[1.0cm]
\end{center}

% \begin{abstract}
% Basic user documentation for GETFEM++.
% \end{abstract}


%%%%%%%%%%%%%%%%%%%%%%%%%%%%%%%%%%%%%%%%%%%%%%%%%%%%%%%%%%%%%%%%%%%%%%%%%
%          INTRODUCTION                                                 %
%%%%%%%%%%%%%%%%%%%%%%%%%%%%%%%%%%%%%%%%%%%%%%%%%%%%%%%%%%%%%%%%%%%%%%%%%

\section*{Introduction}

It is often usefull, when one makes numerical experiments, to compare two solutions defined on the same domain but on different meshes (for instance one mesh more refined than the other) and thus to interpolate one solution defined on a certain mesh to another mesh. GETFEM++ bring two programs located on the directory {\tt contrib} to make those operations in any dimension, provided that the solutions are saved with the format defined by the function {\tt save\_solution} of the file {\tt getfem\_export.h}. This format is easy to exploit and depend on the element (triangle, tetrahedron, quadrilateron ...) and the degree of polynomials used to export the solution. Those programs are independent and can be used in various contexts. The interpolation program allows also to interpolate a N dimension solution on N-1 elements to make a section for instance.\\[5cm]
Copyright (C) 2002\\
The program GETFEM++ is free software; you can redistribute it and/or modify
it under the terms of the GNU General Public License as published by
the Free Software Foundation; version 2 of the License.
This program is distributed in the hope that it will be useful,
but WITHOUT ANY WARRANTY; without even the implied warranty of
MERCHANTABILITY or FITNESS FOR A PARTICULAR PURPOSE.  See the
GNU General Public License for more details.
You should have received a copy of the GNU General Public License
along with this program; if not, write to the Free Software Foundation,
Inc., 59 Temple Place - Suite 330, Boston, MA  02111-1307, USA.

\newpage
\tableofcontents
\newpage

\section{Data format}
The data in the file should begin with a string {\tt DATA ELEMENT}. The text in the file before this string is ignored. Then, the three following lines should be present to define the global caracteristics of the solution
{\tt N = ..\\
  P = ..\\
  K = ..\\
}
where {\tt N} defines the dimension of the domain, {\tt P} defines the dimension of the solution ({\tt P = 1} if the solution is scalar) and {\tt K} defines the degree of polynomials used to save the solution. Just after those three lines and untill the end of the file should be a list of the finite element interpolation nodes and the values of the solution on each node. A blank line should separate the nodes of each element and the list of node is preceded by a line defining the intrinsic dimension of the element (which is less or equal to the dimension of the domain). Then, on each line defining a node, there should be first the coordinates of the node (i.e. {\tt N} values) then the value of each coordinate of the solution (i.e. {\tt P} values). For exemple if a triangle and a tetrahedron are present in a 3D mesh, there would be the following lines (if {\tt N = 3}, {\tt P = 1} and {\tt K = 1})\\[0.5cm]
{\tt 
  DIM = 2 \\
  0 0 0 1.0 \\
  0 1 0 0.5 \\
  0 0 1 0.25 \\
  \\
  DIM = 3 \\
  0 0 0 1.0 \\
  1 0 0 3.0 \\
  0 1 0 0.5 \\
  0 0 1 0.25 \\
  \\
}

\section{Interpolation program}
\section{Comparaison program}



\begin{thebibliography}{99}
% \bibliographystyle{apalike}
% \bibliographystyle{plain}
% \bibliography{all}
\bibitem{USER_DOC} 
  Y. {\sc Renard},
  {\it Getfem++ short user documentation} (2002).
\bibitem{BAS_COMP}
  Y. {\sc Renard},
  {\it Elementary Computations in GETFEM++} (2002).
\end{thebibliography}

\end{document}
